\subsection{The defect cube}\label{sec:construction}

Fix $t\ge1$, $0<\eps<1$ and $0\le r\le(1-\eps)^{t-1}$. The witness is a joint law $\Gamma_t$ on the $3t^2$ copied observations built from independent bits indexed by the cells of the cube $[t]^3$.

\paragraph{Roots.}
Let $D_{ijk}\sim\Bern(\eps)$, $(i,j,k)\in[t]^3$, be mutually independent \emph{defect} bits, and let $N_v\sim\Bern(s)$, one for each copied observation $v$, be mutually independent \emph{private} bits, independent of the defects, with
\begin{equation}\label{eq:s}
s=1-\frac{r}{(1-\eps)^{t-1}}\in[0,1].
\end{equation}
The endpoint $r=0$ uses $s=1$; the endpoint $r=(1-\eps)^{t-1}$ uses $s=0$, in which case the private bits play no role.

\paragraph{Outputs.}
Define
\begin{equation}\label{eq:outputs}
\begin{aligned}
A^{ij}&=\one\bigl\{N_{A^{ij}}=0\text{ and }D_{ijk}=0\ \forall k\bigr\},\\
B^{ik}&=\one\bigl\{N_{B^{ik}}=0\text{ and }D_{ijk}=0\ \forall j\bigr\},\\
C^{jk}&=\one\bigl\{N_{C^{jk}}=0\text{ and }D_{ijk}=0\ \forall i\bigr\}.
\end{aligned}
\end{equation}
So $A^{ij}$ reads the line of cells $\Lambda(A^{ij})=\{(i,j,k):k\in[t]\}$, $B^{ik}$ the line $\Lambda(B^{ik})=\{(i,j,k):j\in[t]\}$, and $C^{jk}$ the line $\Lambda(C^{jk})=\{(i,j,k):i\in[t]\}$; each output equals $1$ exactly when its private bit and every defect on its line are $0$. Let $\Gamma_t$ be the joint law of the outputs. For $s=0$ its entries are
\begin{equation}\label{eq:table}
\Gamma_t(w)=\sum_{d\in\{0,1\}^{[t]^3}:\,F_t(d)=w}\eps^{|d|}(1-\eps)^{t^3-|d|},
\end{equation}
For rational $\eps$, these entries are rational. Here $F_t$ is the map \eqref{eq:outputs} from defect configurations to observed assignments. The entries are nonnegative and sum to $(\eps+1-\eps)^{t^3}=1$.

$\Gamma_t$ does not come from running an inflated triangle model; it is a probability law on the copied observations, which is what a finite inflation test asks for.

\begin{lemma}[Disjoint ancestry gives disjoint inputs]\label{lem:disjoint}
If two sets of copied observations have disjoint copied latent ancestry in the order-$t$ inflation, then their defect lines and their private bits are disjoint, and the two sets of outputs are independent under $\Gamma_t$.
\end{lemma}
\begin{proof}
Two lines through the cube intersect only if they share the index that corresponds to a common source type: $\Lambda(A^{ij})\cap\Lambda(B^{i'k})\ne\emptyset$ forces $i=i'$, and then $X_i$ is a common ancestor; similarly $\Lambda(A^{ij})\cap\Lambda(C^{j'k})\ne\emptyset$ forces $j=j'$ and $\Lambda(B^{ik})\cap\Lambda(C^{jk'})\ne\emptyset$ forces $k=k'$; two lines of the same type are disjoint unless the observations coincide. Taking unions, disjoint ancestry implies disjoint defect supports. Distinct observations have distinct private bits. Outputs that are functions of disjoint families of independent bits are independent.
\end{proof}

\begin{lemma}[Law of a copied triangle]\label{lem:triangle-law}
For every $(i,j,k)\in[t]^3$, $\Law_{\Gamma_t}(A^{ij},B^{ik},C^{jk})=Q(\eps,r)$.
\end{lemma}
\begin{proof}
The three lines of $\Delta_{ijk}$ meet pairwise exactly at the cell $(i,j,k)$. Condition on $D_{ijk}$. If $D_{ijk}=1$, all three outputs are $0$. If $D_{ijk}=0$, each output depends on its own private bit and on the $t-1$ remaining cells of its line; these three families are pairwise disjoint, so the outputs are conditionally independent, each equal to $1$ with probability $(1-s)(1-\eps)^{t-1}=r$ by \eqref{eq:s}. Averaging over $D_{ijk}$ gives \eqref{eq:Q}.
\end{proof}

\begin{lemma}[Symmetry]\label{lem:symmetry}
$\Gamma_t$ is invariant under independent permutations of the $X$-, $Z$- and $Y$-copy indices.
\end{lemma}
\begin{proof}
A permutation $\sigma$ of the $X$-indices sends $D_{ijk}\mapsto D_{\sigma(i)jk}$, $N_{A^{ij}}\mapsto N_{A^{\sigma(i)j}}$, and so on; it permutes the two families of i.i.d.\ bits and carries \eqref{eq:outputs} to itself. The same holds for the other two index families.
\end{proof}
